The following background results are the external dependencies.  The
representation-width argument, recovery by independent tensor operations,
change of quadratic form, and specialization of the prescribed class are
proved in the text.  No algebraicity theorem for the Kuga--Satake class is
an input.  Page numbers for \cite{Huybrechts} refer to the linked draft.

\begingroup
\small
\setlength{\tabcolsep}{5pt}
\renewcommand{\arraystretch}{1.13}
\begin{longtable}{@{}>{\raggedright\arraybackslash}p{.40\textwidth}>{\raggedright\arraybackslash}p{.31\textwidth}>{\raggedright\arraybackslash}p{.23\textwidth}@{}}
\toprule
Established statement and hypotheses & Reference and locator & Role \\
\midrule
\endfirsthead
\toprule
Established statement and hypotheses & Reference and locator & Role \\
\midrule
\endhead
\bottomrule
\endfoot

For smooth projective complex varieties, Hodge decomposition, Poincar\'e
duality, the Lefschetz theorem on $(1,1)$-classes, hard Lefschetz, and the
Hodge--Riemann relations hold; pullback, Gysin maps, and products have their
usual Hodge types. &
\cite{Voisin}, \S\S6.1--6.3, pp.~139--154;
\S\S7.1--7.3, pp.~157--182. &
Pairings, primitive cohomology, divisor classes, and nonvanishing
calculations. \\[3pt]

Polarizable rational Hodge structures have Hodge complements, and their
rational Hodge group is reductive.  The same
holds for a finite joint collection. &
\cite{Deligne}, Proposition 3.6 and its proof, electronic p.~27;
\cite{Voisin}, \S7.3.1. &
Rational splittings and the joint Hodge-group Lie algebra. \\[3pt]

In characteristic zero algebraic group homomorphisms have closed image
and surjective differential onto its Lie algebra.  A unipotent group has
nonzero invariants in every nonzero representation.  A reductive Lie
algebra is the sum of its center and simple ideals. &
\cite{Borel}, Corollary 1.4, Theorem 4.8, \S7,
and Proposition 14.2;
\cite{FultonHarris}, Lectures 9 and 14. &
The rational simple ideal above the orthogonal quotient and the
commuting kernel. \\[3pt]

A complex K3 surface has the rank-$22$ cup lattice and Hodge numbers used
here.  Its universal deformation is smooth of dimension $20$; the local
period map is an isomorphism.  Fixing a nonzero integral $(1,1)$-class
cuts out the corresponding local period hypersurface and extends its
line bundle. &
\cite{Huybrechts}, Chapter 1, \S\S2--3;
Chapter 6, Corollary 2.7 and Proposition 2.8, p.~107;
\S2.4, p.~108. &
Marked polarized deformation charts and the rank-$21$ primitive space. \\[3pt]

The full even Clifford construction of a polarized Hodge structure of K3
type is polarizable of weight one.  Clifford multiplication with an
invertible right multiplier gives the standard Hodge embedding. &
\cite{Huybrechts}, Chapter 4, Propositions 2.5--2.6 and Remark 2.7,
pp.~66--68;
\cite{vanGeemen}, Propositions 5.9 and 6.3. &
The prescribed full target, Hodge types, and fixed construction data. \\[3pt]

Complex semisimple Lie algebras have completely reducible finite-dimensional
representations and highest-weight classification.  In type $B_n$ the
dominant-weight coordinates and spin weights have the stated forms.
The Clifford Basis Theorem and spin double covering hold in characteristic
zero. &
\cite{FultonHarris}, Lectures 14 and 18--20;
\S19.4, pp.~294--296;
\S\S20.1--20.2, pp.~299--312. &
The width bound, spin multiplicity calculation, and finite-kernel lift
comparison. \\[3pt]

The functor $H^1$ identifies complex abelian varieties up to isogeny
contravariantly with polarizable rational weight-one Hodge structures.
A polarized complex torus is projective. &
\cite{Milne}, Theorems 6.7--6.8 and Corollary 6.9,
pp.~72--73 (stated covariantly for homology). &
Algebraic realization of every rational weight-one slot map and of the
Clifford spaces. \\[3pt]

For a smooth projective complex variety $Y$, its Albanese is an abelian
variety, and its Albanese map induces an isomorphism on holomorphic
one-forms. &
\cite{Albanese}, Chapter VI, \S9.2, pp.~315--316,
formulas (9.11)--(9.13). &
Conversion of the surface-cover input into a correspondence involving an
abelian variety. \\[3pt]

On smooth projective varieties, correspondences compose by
pullback, intersection, and proper pushforward, with the indicated
codimension.  Cycle classes commute with these operations. &
\cite{Fulton}, \S16.1, pp.~305--315, and
\S\S19.1--19.2, pp.~371--385;
\cite{StacksCorrespondences}, Lemmas 45.3.1 and 45.3.6. &
Actual rational cycles throughout, with controlled degrees and graph
transposes. \\[3pt]

For a projective family over a noetherian base, each fixed-polynomial
relative Hilbert scheme is projective and has a flat universal
subscheme.  One uniform sufficiently high degree gives its Grassmannian
realization. &
\cite{GrothendieckHilbert}, Theorem 3.2, p.~260;
Proposition 3.8, p.~265. &
Countably many proper cycle parameter spaces and the analytic
classifying map. \\[3pt]

Flat morphisms of finite presentation with smooth fibers are smooth.
In a proper flat family the smooth-fiber locus is open: remove the
proper image of the closed nonsmooth locus.  High powers of an ample
line bundle on a projective variety are very ample and have vanishing
higher cohomology.  Riemann--Roch expresses their Euler characteristics
by Chern classes. &
\cite{StacksFamilies}, tags 01V4, 01V8, 01VU, 0B5U;
\cite{Fulton}, \S18.3. &
The smooth ambient Hilbert family, projective embeddings, and the fixed
Hilbert polynomial. \\[3pt]

If a Cartier equation is a nonzerodivisor on a closed subscheme $Z$ of
dimension at most $k+1$, its Cartier pullback takes $[Z]_{k+1}$ to the
$k$-dimensional cycle of the cut, with scheme multiplicities. &
\cite{StacksCartier}, Lemma 42.29.5(1), tag 02TB;
\cite{Fulton}, Chapters 2 and 6. &
Specialization of top cycles, including nonreduced and nonpure Hilbert
fibers. \\[3pt]

A proper holomorphic map has coherent higher direct images.  For a
coherent sheaf flat over a reduced base, fiber cohomology is upper
semicontinuous; constant dimensions give local freeness and fiber base
change.  Coherent sheaves on Stein spaces have vanishing higher
cohomology. &
\cite{Grauert}, \S6, Hauptsatz I, p.~59;
\S7, S\"atze 3 and 5, pp.~62, 64;
\S2, Satz 3, p.~17, with the 1963 correction. &
Extension of polarizations, local relative projective embeddings, and
holomorphic Grassmannian maps. \\[3pt]

A proper smooth submersion is locally differentiably trivial.  Proper
holomorphic images of closed analytic sets are analytic.  Irreducible
components of a complex analytic set are closed and locally finite;
in a second-countable space there are only countably many. &
\cite{Ehresmann}, Lemma 9.2, p.~77;
\cite{Grauert}, \S7, Satz 1, p.~60;
\cite{Albanese}, Chapter II, Theorem 5.3, pp.~102--103. &
Marked flat cohomology, proper-image avoidance, and transport to the
chosen fiber. \\[3pt]

Complex quasi-projective varieties admit projective resolutions.
Closed analytic subvarieties of projective space and holomorphic maps
between projective varieties are algebraic.  Projective coherent
cohomology agrees with analytic cohomology. &
\cite{Hironaka}, Part I, Main Theorems I and I$^*$;
\cite{SerreGAGA}, Theorems 1--3, pp.~19--20,
and Propositions 13--15, pp.~29--30. &
Smooth cycle parameters and surfaces; algebraicity of the resulting
fiber cycles and weight-one homomorphisms. \\

\end{longtable}
\endgroup
